\documentclass[11pt]{article}
\usepackage[margin=0.9in]{geometry}
\usepackage[T1]{fontenc}
\usepackage{lmodern}
\usepackage{amsmath,amssymb,amsthm}
\usepackage[hidelinks]{hyperref}
\usepackage{microtype}
\usepackage{xcolor}
\usepackage{fancyhdr}

\definecolor{accent}{RGB}{34,78,112}
\newtheorem*{conclusion}{Conclusion}
\pagestyle{fancy}
\fancyhf{}
\setlength{\headheight}{14pt}
\lhead{Literature-status note: arXiv:1103.0223}
\rhead{Lane 4}
\cfoot{\thepage}
\fancypagestyle{plain}{%
  \fancyhf{}%
  \lhead{Literature-status note: arXiv:1103.0223}%
  \rhead{Lane 4}%
  \cfoot{\thepage}%
}
\setlength{\parindent}{0pt}
\setlength{\parskip}{0.55em}

\begin{document}

\begin{center}
  {\Large\bfseries\color{accent} Polynomial ergodic averages on arbitrary intervals}\\[4pt]
  {\large A literature-implied answer to Austin's first further question}\\[6pt]
  {\small Status: \texttt{literature\_implied\_answer} \quad Updated 27 August 2026}
\end{center}

\paragraph{Source.}
Tim Austin, \emph{Norm convergence of continuous-time polynomial multiple
ergodic averages}, arXiv:1103.0223.  In Section~4, ``Further questions'',
PDF page~26, Austin asks whether Theorem~1.2 remains valid over arbitrary real
intervals whose lengths tend to infinity, independently of their distance from
the origin.  In the paper's notation, for a measurable measure-preserving
$\mathbb R^D$-action $\tau$, polynomial maps $p_i:\mathbb R\to\mathbb R^D$,
and $f_i\in L^\infty(\mu)$, the question is whether
\[
  A_I(f_1,\ldots,f_k)
  :=\frac{1}{|I|}\int_I \prod_{i=1}^k f_i\circ\tau^{p_i(t)}\,dt
\]
converges in $L^2(\mu)$ whenever $|I|\to\infty$, with no restriction on the
location of $I$.

\paragraph{Supporting theorem.}
Pavel Zorin-Kranich, \emph{Norm convergence of multiple ergodic averages on
amenable groups}, arXiv:1111.7292, Theorem~1.1(1), PDF page~2, proves the
following.  Let $\mathcal G$ be a locally compact amenable group,
$\mathcal T$ a nilpotent group of measure-preserving Koopman operators, and
$T_i:\mathcal G\to\mathcal T$ measurable polynomial maps.  Then the associated
multiple averages converge in $L^2$ along every left Folner net, and the limit
does not depend on the Folner net.  Section~4 of that paper defines polynomial
maps by iterated discrete derivatives; its example following Theorem~4.5
explicitly includes conventional real polynomials composed with one-parameter
subgroups.

\paragraph{Exact identification.}
Set $\mathcal G=(\mathbb R,+)$ and let $U_v f=f\circ\tau^v$.  The image
\[
  \mathcal T=\{U_v:v\in\mathbb R^D\}
\]
is abelian, hence nilpotent.  Each $T_i(t):=U_{p_i(t)}$ is a measurable
polynomial map in Zorin-Kranich's sense.  Thus Theorem~1.1 applies to Austin's
integrand (take the harmless factor $f_0=1$).

Now let $I_n$ be arbitrary intervals with $|I_n|\to\infty$.  For every compact
$K\subset\mathbb R$,
\[
 \sup_{h\in K}\frac{|(h+I_n)\mathbin{\triangle}I_n|}{|I_n|}
 \leq \frac{2\sup_{h\in K}|h|}{|I_n|}\longrightarrow 0.
\]
Hence $(I_n)$ is a Folner sequence, regardless of how far its members lie from
the origin.  Zorin-Kranich's theorem therefore gives the requested $L^2$
convergence and shows that the limit is independent of the chosen moving
interval sequence.

\begin{conclusion}
Austin's arbitrary-interval question has an affirmative answer: the averages
$A_I(f_1,\ldots,f_k)$ converge in $L^2(\mu)$ as $|I|\to\infty$, uniformly in
the sense that every sequence of intervals with diverging lengths has the same
limit.
\end{conclusion}

\paragraph{Additional scope.}
The same specialization with $\mathcal G=\mathbb R^r$ also answers Austin's
Question~4.1 (PDF page~26): polynomial averages over $[0,R)^r$ converge, and in
fact they converge along arbitrary Folner nets in $\mathbb R^r$.  It also
covers the final extension on PDF page~27 to actions of a connected nilpotent
Lie group, provided ``polynomial'' is understood in the discrete-derivative
(Leibman) sense used by the supporting theorem.  The intervening Hardy-field
question is outside this polynomial theorem and is not classified here.

\paragraph{Provenance and novelty status.}
This is not a new theorem produced by the run.  Zorin-Kranich cites Austin's
work but does not explicitly identify the Section~4 question as the target of
Theorem~1.1; the relation above is a direct specialization found during this
run.  That is why the packet is classified as a literature-implied answer
rather than an explicitly literature-answered problem.  A bounded search by
arXiv id, exact title, ``arbitrary intervals'', ``uniform Cesaro'', and
``Folner'' also found Bergelson--Leibman--Moreira's later Folner-sequence
theorem for one-parameter flows; Zorin-Kranich's theorem supplies the decisive
full $\mathbb R^D$-action scope.

\paragraph{Human review recommendation.}
High confidence.  Verify only the standard functorial point that composing the
$\mathbb R^D$-action with an ordinary vector polynomial gives a polynomial
operator-valued map in Definition~4.3 of arXiv:1111.7292.  The abelian case is
also illustrated explicitly after its Theorem~4.5.

\begin{thebibliography}{9}
\bibitem{Austin}
T. Austin, \emph{Norm convergence of continuous-time polynomial multiple
ergodic averages}, Ergodic Theory Dynam. Systems 32 (2012), 361--382;
arXiv:1103.0223.

\bibitem{ZK}
P. Zorin-Kranich, \emph{Norm convergence of multiple ergodic averages on
amenable groups}, J. Analyse Math. 130 (2016), 219--241;
arXiv:1111.7292.

\bibitem{BLM}
V. Bergelson, A. Leibman, and C. G. Moreira,
\emph{From discrete- to continuous-time ergodic theorems},
Ergodic Theory Dynam. Systems 32 (2012), 383--426.
\end{thebibliography}

\end{document}
